% ===== 原文 pp.47–61（PDF p.62–76），tr-03：第 2 章 §2.1 后段 =====

% [原文 p.47]
\begin{theorem}
第一结构联络是平坦的，当且仅当 $\cM$ 为 Frobenius 流形，且 $E$ 是满足
\[
[E,X\circ Y]-[E,X]\circ Y-X\circ[E,Y]=d_0X\circ Y
\]
的 Euler 向量场。
\end{theorem}

证明参见 [76] 的定理 I.2.5.2。

\subsubsection{量子微分方程}

假设给定一个非异簇 $X$，其上同调只出现在偶次。照例取 $T_0,\dots,T_m$ 为 $H^*(X,\C)$ 的一组生成元，其中 $T_0$ 是 $X$ 的基本类，$T_1,\dots,T_p$ 是 $H^2(X,\R)$ 的一组生成元，且在所有由稳定映射代表的 $\beta\in H_2(X,\Z)$ 上取值非负。设 $T^0,\dots,T^m$ 为 Poincaré 对偶基。如在例 2.8 中那样，这给出一个 Frobenius 流形 $\cM$，带有 Euler 向量场 $E$ 与第一结构联络 $\widehat{\nabla}$。下面我们描述\emph{量子微分方程}
\begin{equation}
\widehat{\nabla}_{\partial_{y_i}}s=0,\qquad i=0,\dots,m. \tag{2.8}
\end{equation}
的解。特别地，我们将写下该方程的一组基本解。

我们首先需要

\begin{proposition}
拓扑递推关系（TRR）。
\begin{align*}
&\langle\psi^{d_1}\alpha_1,\dots,\psi^{d_n}\alpha_n\rangle_{0,\beta}\\
&\quad=\sum\langle\psi^{d_1-1}\alpha_1,\prod_{i\in S_1}\psi^{d_i}\alpha_i,T_e\rangle_{0,\beta_1}\langle T^e,\psi^{d_2}\alpha_2,\psi^{d_3}\alpha_3,\prod_{i\in S_2}\psi^{d_i}\alpha_i\rangle_{0,\beta_2}
\end{align*}
其中求和遍历所有 $0\le e\le m$、所有分解 $\beta_1+\beta_2=\beta$，以及满足 $S_1\cup S_2=\{4,\dots,n\}$ 的不交指标集 $S_1,S_2$。
\end{proposition}

证明可参见例如 [90]。

\begin{remark}
事实上，拓扑递推关系表明，人们可以由通常的亏格零 Gromov–Witten 不变量完全重构亏格零后代不变量。因此亏格零后代 Gromov–Witten 不变量实际上并不携带新信息；但在这里记录它们是有用的，其确切原因正是它们给出量子微分方程的解。
\end{remark}

% [Tutor Note: 命题 2.12 的求和号在原文中没有写指标集，条件“$0\le e\le m$、$\beta_1+\beta_2=\beta$、$S_1\cup S_2=\{4,\dots,n\}$”以文字给出（见 p.47）；注意 $S_1,S_2$ 只划分 $\{4,\dots,n\}$，前三个插入点 $\alpha_1,\alpha_2,\alpha_3$ 被单独处理，这是 TRR 的标准形式。]

下面我们引入另一种记号：对 $\alpha_1,\dots,\alpha_n\in H^*(X,\C)$ 与非负整数 $d_1,\dots,d_n$，记
\[
\langle\langle\psi^{d_1}\alpha_1,\dots,\psi^{d_n}\alpha_n\rangle\rangle
:=\sum_{k=0}^{\infty}\sum_{\beta\in H_2(X,\Z)}\frac{1}{k!}\langle\psi^{d_1}\alpha_1,\dots,\psi^{d_n}\alpha_n,\gamma^k\rangle_{0,\beta}.
\]
这里，如同在 Gromov–Witten 势的定义中那样，$\gamma=\sum_{i=0}^m y_iT_i$，而 $\gamma^k$ 表示把 $\gamma$ 取 $k$ 次。注意这让我们可以把 Gromov–Witten 势写成
\[
\Phi=\langle\langle1\rangle\rangle.
\]
此外，
\[
\partial_{y_i}\langle\langle\psi^{d_1}\alpha_1,\dots,\psi^{d_n}\alpha_n\rangle\rangle
=\langle\langle\psi^{d_1}\alpha_1,\dots,\psi^{d_n}\alpha_n,T_i\rangle\rangle.
\]
于是对 $0\le i\le m$ 可定义
\[
s_i:=T_i-\sum_{j=0}^m\Big\langle\Big\langle\frac{T_i}{\hbar+\psi},T_j\Big\rangle\Big\rangle T^j.
\]

% [原文 p.48]
这里 $T_i/(\hbar+\psi)$ 按如下幂级数展开来理解：
\[
\frac{T_i}{\hbar+\psi}=\sum_{n=0}^{\infty}(-1)^n\hbar^{-(n+1)}\psi^nT_i.
\]

\begin{theorem}
诸 $s_i$ 满足量子微分方程 (2.8)。
\end{theorem}

\begin{proof}[证明]
TRR 给出
\begin{align}
&\langle\langle\psi^{d_1+1}T_{j_1},\psi^{d_2}T_{j_2},\psi^{d_3}T_{j_3}\rangle\rangle\notag\\
&\quad=\sum_{k=0}^{\infty}\sum_{\beta}\frac{1}{k!}\langle\psi^{d_1+1}T_{j_1},\psi^{d_2}T_{j_2},\psi^{d_3}T_{j_3},\gamma^k\rangle_{0,\beta}\notag\\
&\quad=\sum_{k=0}^{\infty}\sum_{i=0}^m\sum_{\substack{\beta_1,\beta_2\\S_1,S_2}}\frac{1}{k!}\langle\psi^{d_1}T_{j_1},\gamma^{\#S_1},T_i\rangle_{0,\beta_1}\langle T^i,\psi^{d_2}T_{j_2},\psi^{d_3}T_{j_3},\gamma^{\#S_2}\rangle_{0,\beta_2}\tag{2.9}\\
&\quad=\sum_{k=0}^{\infty}\sum_{i=0}^m\sum_{\substack{\beta_1,\beta_2\\a+b=k}}\frac{1}{a!b!}\langle\psi^{d_1}T_{j_1},\gamma^a,T_i\rangle_{0,\beta_1}\langle T^i,\psi^{d_2}T_{j_2},\psi^{d_3}T_{j_3},\gamma^b\rangle_{0,\beta_2}\notag\\
&\quad=\sum_{i=0}^m\langle\langle\psi^{d_1}T_{j_1},T_i\rangle\rangle\langle\langle T^i,\psi^{d_2}T_{j_2},\psi^{d_3}T_{j_3}\rangle\rangle.\notag
\end{align}
注意量子微分方程 (2.8) 可对 $s=s_j$ 写成
\begin{equation}
\hbar\partial_{y_i}(s_j)=-T_i*s_j,\qquad\text{对 }0\le i\le m. \tag{2.10}
\end{equation}
现在
\[
\hbar\partial_{y_i}(s_j)=-\sum_{k=0}^m\Big\langle\Big\langle T_i,\frac{\hbar T_j}{\hbar+\psi},T_k\Big\rangle\Big\rangle T^k.
\]
另一方面，$\partial_{y_i}\partial_{y_j}\partial_{y_k}\Phi=\langle\langle T_i,T_j,T_k\rangle\rangle$，而量子乘积由
\[
T_i*T_j=\sum_k\langle\langle T_i,T_j,T_k\rangle\rangle T^k
\]
给出。于是
\begin{align*}
T_i*s_j &= T_i*T_j-\sum_{k=0}^m\Big\langle\Big\langle\frac{T_j}{\hbar+\psi},T_k\Big\rangle\Big\rangle T_i*T^k\\
&= \sum_{\ell=0}^m\langle\langle T_i,T_j,T_\ell\rangle\rangle T^\ell-\sum_{k,\ell=0}^m\Big\langle\Big\langle\frac{T_j}{\hbar+\psi},T_k\Big\rangle\Big\rangle\langle\langle T_i,T^k,T_\ell\rangle\rangle T^\ell\\
&= \sum_{\ell=0}^m\langle\langle T_i,T_j,T_\ell\rangle\rangle T^\ell-\sum_{\ell=0}^m\Big\langle\Big\langle\frac{\psi T_j}{\hbar+\psi},T_i,T_\ell\Big\rangle\Big\rangle T^\ell\\
&= \sum_{\ell=0}^m\Big\langle\Big\langle T_i,\frac{\hbar T_j}{\hbar+\psi},T_\ell\Big\rangle\Big\rangle T^\ell.
\end{align*}
这里倒数第二个等号由式 (2.9) 得到。因此 (2.10) 的左、右两端一致。
\end{proof}

% [Tutor Note: 式 (2.9) 在原文里是五行连等式，编号只落在第 3 行（见 p.48 图），译文的编号 2.9 也放在第 3 行以免错位；核对时注意第 3、4 行的两个求和都从 $k=0$ 起，第 4 行把 $k$ 按 $a+b=k$ 拆分，这一步正是 $\langle\langle\cdot\rangle\rangle$ 的 $\frac{1}{k!}$ 与两个 $\frac{1}{a!b!}$ 的配合处。]

利用后代 Gromov–Witten 不变量的公理把诸 $s_i$ 改写如下，将是有用的。

% [原文 p.49]
\begin{proposition}
\begin{align*}
s_i &= e^{-(\sum_{k=0}^p y_kT_k)/\hbar}\cup T_i-\sum_{\substack{\beta,j\\n_k}}\Big\langle\frac{e^{-(\sum_{k=0}^p y_kT_k)/\hbar}\cup T_i}{\hbar+\psi},T_j,T_{p+1}^{n_{p+1}},\dots,T_m^{n_m}\Big\rangle_{0,\beta}T^j\\
&\qquad\qquad\cdot\Big(\prod_{\ell=1}^p e^{y_\ell\int_\beta T_\ell}\Big)\frac{y_{p+1}^{n_{p+1}}\cdots y_m^{n_m}}{n_{p+1}!\cdots n_m!}
\end{align*}
其中求和遍历 $0\le j\le m$、$\beta\in H_2(X,\Z)$ 以及 $n_{p+1},\dots,n_m\ge0$。表达式 $e^{-(\sum_{i=0}^p y_iT_i)/\hbar}$ 在环 $H^*(X,\C)\otimes_\C\C[y_1,\dots,y_p][\![y_0,\hbar^{-1}]\!]$ 中解释，配备通常的杯积。
\end{proposition}

% [Tutor Note: 命题 2.15 的含义是：$s_i$ 的“指数因子”把 $y_0,\dots,y_p$ 合并到 $\hbar^{-1}$ 的指数上，而 $y_{p+1},\dots,y_m$（非除子类方向的参数）只以普通幂次出现。这一分界来自除子公理：$T_1,\dots,T_p$ 是除子类，故 $\exp(y_\ell\int_\beta T_\ell)$ 因子可以被提出；核对锚点是环 $\C[y_1,\dots,y_p][\![y_0,\hbar^{-1}]\!]$ 中的双括号——$y_0$ 与 $\hbar^{-1}$ 是形式变量。]

\begin{proof}[证明]
首先注意，由基本类公理，再由后代不变量的除子公理，若 $\beta\ne0$ 或 $\beta=0$ 且 $\sum_{k=p+1}^m n_k\ge1$，则
\begin{align*}
&\Big\langle\frac{T_i}{\hbar+\psi},T_j,T_0^{n_0},\dots,T_m^{n_m}\Big\rangle_{0,\beta}\\
&\quad=\Big\langle\sum_{n=0}^{\infty}(-1)^n\hbar^{-(n+1)}\psi^nT_i,T_j,T_0^{n_0},\dots,T_m^{n_m}\Big\rangle_{0,\beta}\\
&\quad=\Big\langle\sum_{n=n_0}^{\infty}(-1)^n\hbar^{-(n+1)}\psi^{n-n_0}T_i,T_j,T_1^{n_1},\dots,T_m^{n_m}\Big\rangle_{0,\beta}\\
&\quad=\sum_{n_k'}\Big\langle\sum_n(-1)^n\hbar^{-(n+1)}\psi^{n-n_0-n_1'-\cdots-n_p'}(T_i\cup T_1^{n_1'}\cup\cdots\cup T_p^{n_p'}),T_j,T_{p+1}^{n_{p+1}},\dots,T_m^{n_m}\Big\rangle_{0,\beta}\\
&\qquad\qquad\qquad\qquad\cdot\prod_{\ell=1}^p\binom{n_\ell}{n_\ell'}\Big(\int_\beta T_\ell\Big)^{n_\ell-n_\ell'}
\end{align*}
其中对 $n_k'$ 的求和遍历所有满足 $0\le n_k'\le n_k$ 的 $n_1',\dots,n_p'$，而对 $n$ 的求和遍历 $n\ge n_0+n_1'+\cdots+n_p'$。又注意，当 $\beta=0$ 时，若 $k>p$ 时 $n_k=0$，但某个 $0\le k\le p$ 的 $n_k\ne0$，则由基本类公理、除子公理与点映射公理，
\begin{align*}
&\Big\langle\frac{T_i}{\hbar+\psi},T_j,T_0^{n_0},\dots,T_p^{n_p}\Big\rangle_{0,0}\\
&\quad=\langle(-1)^{n_0+\cdots+n_p-1}\hbar^{-(n_0+\cdots+n_p)}(T_i\cup T_0^{n_0}\cup\cdots\cup T_k^{n_k-1}\cup\cdots\cup T_p^{n_p}),T_j,T_k\rangle_{0,0}\\
&\quad=(-1)^{n_0+\cdots+n_p-1}\hbar^{-(n_0+\cdots+n_p)}\int_X T_i\cup T_j\cup T_0^{n_0}\cup\cdots\cup T_p^{n_p}.
\end{align*}
于是，若把 $\langle\langle T_i/(\hbar+\psi),T_j\rangle\rangle$ 按 $\gamma^k$ 展开，我们便把它分成两类项：涉及 $\beta=0$ 且不含 $k\ge p+1$ 的 $T_k$ 的项，以及另一类项。
% [原文 p.50]
这给出
\begin{align*}
&\Big\langle\Big\langle\frac{T_i}{\hbar+\psi},T_j\Big\rangle\Big\rangle\\
&\quad=\sum_{n_0+\cdots+n_p\ge1}(-1)^{n_0+\cdots+n_p-1}\hbar^{-(n_0+\cdots+n_p)}\int_X T_i\cup T_j\cup T_0^{n_0}\cup\cdots\cup T_p^{n_p}\frac{y_0^{n_0}\cdots y_p^{n_p}}{n_0!\cdots n_p!}\\
&\quad\quad+\sum_{\beta,n_k}\sum_{n_k'}\Big\langle\sum_n(-1)^n\hbar^{-(n+1)}\psi^{n-n_0-n_1'-\cdots-n_p'}(T_i\cup T_1^{n_1'}\cup\cdots\cup T_p^{n_p'}),T_j,T_{p+1}^{n_{p+1}},\dots,T_m^{n_m}\Big\rangle_{0,\beta}\\
&\qquad\qquad\qquad\qquad\cdot\Big(\prod_{\ell=1}^p\binom{n_\ell}{n_\ell'}\Big(\int_\beta T_\ell\Big)^{n_\ell-n_\ell'}\Big)\frac{y_0^{n_0}\cdots y_m^{n_m}}{n_0!\cdots n_m!},
\end{align*}
其中第二项中的求和遍历所有满足“并非同时有 $\beta=0$ 与 $n_{p+1}=\cdots=n_m=0$”的 $\beta,n_k$，以及 $0\le n_k'\le n_k$。这于是等同于
\begin{align*}
&\sum_{n_0+\cdots+n_p\ge1}(-1)^{n_0+\cdots+n_p-1}\hbar^{-(n_0+\cdots+n_p)}\int_X T_i\cup T_j\cup T_0^{n_0}\cup\cdots\cup T_p^{n_p}\frac{y_0^{n_0}\cdots y_p^{n_p}}{n_0!\cdots n_p!}\\
&\quad+\sum_{\beta,n_k}\sum_{n_k'}\sum_{n=0}^{\infty}\Big\langle\frac{(-1)^n\hbar^{-(n+1)}\psi^n(T_i\cup(y_0T_0)^{n_0}\cup\bigcup_{\ell=1}^p(y_\ell T_\ell)^{n_\ell'})(-\hbar)^{-(n_0+n_1'+\cdots+n_p')}}{n_0!(n_1')!\cdots(n_p')!},\\
&\qquad\qquad\qquad\qquad T_j,T_{p+1}^{n_{p+1}},\dots,T_m^{n_m}\Big\rangle_{0,\beta}\prod_{\ell=1}^p\frac{1}{(n_\ell-n_\ell')!}\Big(y_\ell\int_\beta T_\ell\Big)^{n_\ell-n_\ell'}\frac{y_{p+1}^{n_{p+1}}\cdots y_m^{n_m}}{n_{p+1}!\cdots n_m!}\\
&\quad=-\int_X T_i\cup T_j\cup(e^{-(y_0T_0+\cdots+y_pT_p)/\hbar}-1)\\
&\quad\quad+\sum_{\beta,n_k}\Big\langle\frac{T_i\cup e^{-(y_0T_0+\cdots+y_pT_p)/\hbar}}{\hbar+\psi},T_j,T_{p+1}^{n_{p+1}},\dots,T_m^{n_m}\Big\rangle_{0,\beta}\\
&\qquad\qquad\qquad\qquad\cdot\Big(\prod_{\ell=1}^p e^{y_\ell\int_\beta T_\ell}\Big)\frac{y_{p+1}^{n_{p+1}}\cdots y_m^{n_m}}{n_{p+1}!\cdots n_m!}.
\end{align*}
这便给出 $s_i$ 所需的形式。
\end{proof}

我们仍希望理解两个问题。第一，按我们的描述，量子微分方程的解 $s_i$ 纯属形式解，我们希望理解它们是否收敛。第二，我们已经证明诸 $s_i$ 在水平（$\cM$）方向上关于联络 $\widehat{\nabla}$ 是 $\cM\times\C^\times$ 上的平坦截面，但我们想理解它们在竖直（$\C^\times$）方向上的行为。

这两个问题密切相关，但首先我们将引入一个替代 $\cM$ 的对象，从今往后讨论量子上同调时都在它上面工作。

回忆 $\cM=\mathrm{Spec}\,\C[\kappa_1^{\pm},\dots,\kappa_p^{\pm}][\![y_0,y_{p+1},\dots,y_m]\!]$。我们宁愿把它看成某个复解析空间的一种形式完备化。但由于今后事实上必须在 $\cM$ 的万有覆叠 $\widetilde{\cM}$ 上工作，我们将直接描述 $\widetilde{\cM}$。

% [原文 p.51]
\begin{definition}
我们定义 $\widetilde{\cM}$ 为环空间 $(\C^p,\cO_{\widetilde{\cM}})$，其中 $\C^p$ 上的坐标为 $y_1,\dots,y_p$，而 $\cO_{\widetilde{\cM}}$ 表示 $\widetilde{\cM}$ 上的结构层：在开集 $U\subseteq\widetilde{\cM}$ 上，$\cO_{\widetilde{\cM}}(U)$ 由形式幂级数
\begin{equation}
\sum f_{i_0i_{p+1}\cdots i_m}y_0^{i_0}y_{p+1}^{i_{p+1}}\cdots y_m^{i_m} \tag{2.11}
\end{equation}
组成，其中 $f_{i_0i_{p+1}\cdots i_m}$ 是开集 $U$ 上的全纯函数。
\end{definition}

有一个映射 $\widetilde{\cM}\to\cM$，由 $\kappa_i=e^{y_i}$（$1\le i\le p$）与 $y_j=y_j$（$j\notin\{1,\dots,p\}$）给出。从今往后，当我们讨论由量子上同调产生的结构时，都在这个形式设定下工作。

我们还定义一些额外的记号。

% [Tutor Note: 定义 2.16 中 $\cM$ 的坐标（见 p.50 末）是 $\C[\kappa_1^{\pm},\dots,\kappa_p^{\pm}][\![y_0,y_{p+1},\dots,y_m]\!]$，而 $\widetilde{\cM}$ 用 $y_1,\dots,y_p$ 代替 $\kappa_1,\dots,\kappa_p$：$\kappa_i=e^{y_i}$ 只在 $i=1,\dots,p$ 上取指数，$y_0,y_{p+1},\dots,y_m$ 原样保留。这正是“复解析空间的形式完备化”的确切含义，也是下文所有 $\cO_{\widetilde{\cM}}\{\hbar,\hbar^{-1}\}$ 记号的基础。]

\begin{definition}
记 $\C\{\hbar,\hbar^{-1}\}$ 为在某个去心圆盘 $\{\hbar\in\C\mid0<|\hbar|<\epsilon\}$ 上收敛的 Laurent 级数之环（$\epsilon$ 取某个值）。设 $\C\{\hbar\}$ 为在 $\hbar=0$ 处全纯的函数构成的子环，并设 $\cO(\PP^1\setminus\{0\})$ 为 $\PP^1\setminus\{0\}$ 上的全纯函数之环（坐标为 $\hbar$）。我们可以写
\[
\C\{\hbar,\hbar^{-1}\}=\C\{\hbar\}\oplus\hbar^{-1}\cO(\PP^1\setminus\{0\}).
\]
下面我们定义
\[
\cO_{\widetilde{\cM}}\{\hbar,\hbar^{-1}\}
\]
为这样一个层：它在开集 $U$ 上的取值由如式 (2.11) 的形式级数组成，使得每个系数 $f_{i_0i_{p+1}\cdots i_m}$ 都是
\[
\{(y,\hbar)\in U\times\C\mid0<|\hbar|<\epsilon(y)\}
\]
上的全纯函数，其中 $\epsilon:U\to\R_{>0}$ 为某个连续映射。于是可定义
\[
\cO_{\widetilde{\cM}}\{\hbar\}\subseteq\cO_{\widetilde{\cM}}\{\hbar,\hbar^{-1}\}
\]
为系数 $f_{i_0i_{p+1}\cdots i_m}$ 在 $\hbar=0$ 处全纯的子层。类似地可用同样的方式定义 $\cO_\cM\{\hbar\}$ 与 $\cO_\cM\{\hbar,\hbar^{-1}\}$。
\end{definition}

我们将常常希望考虑 $\cM$ 只是通常的复流形的情形，因为这在概念上稍微容易一些。此时 $\cO_\cM\{\hbar,\hbar^{-1}\}$ 是这样的层：它在开集 $U\subseteq\cM$ 上由在 $\{(y,\hbar)\in U\times\C\mid0<\hbar<\epsilon(y)\}$ 上全纯的函数 $f$ 组成，其中 $\epsilon:U\to\R_{>0}$ 为某个连续映射（依赖于 $f$）。此外，$\cO_\cM\{\hbar\}$ 是由在 $\hbar=0$ 处也全纯的函数构成的子层。

回到关于 $s_i$ 的收敛性以及 $s_i$ 在 $\hbar$ 方向上行为的问题，让我们引入映射 $S$，它把
\[
\alpha\in H^*(X,\C)\otimes_\C\cO_{\widetilde{\cM}}\{\hbar,\hbar^{-1}\}
\]
送到
\[
S(\alpha)=\alpha-\sum_{j=0}^m\Big\langle\Big\langle\frac{\alpha}{\hbar+\psi},T_j\Big\rangle\Big\rangle T^j.
\]
我们希望把 $S$ 看成映射
\[
S:H^*(X,\C)\otimes_\C\cO_{\widetilde{\cM}}\{\hbar,\hbar^{-1}\}\to H^*(X,\C)\otimes_\C\cO_{\widetilde{\cM}}\{\hbar,\hbar^{-1}\},
\]
但先验地，$S(\alpha)$ 只是纯形式的表达式，所以此处我们暂不描述 $S$ 的像。不过，我们可以用 $S$ 来计算 $\widehat{\nabla}_{\hbar\partial_\hbar}$ 对 $s_i$ 的作用：

% [原文 p.52]
\begin{proposition}
\[
\widehat{\nabla}_{\hbar\partial_\hbar}s_i=S(\mathrm{Gr}_E(T_i)-\hbar^{-1}c_1(\cT_X)\cup T_i).
\]
\end{proposition}

\begin{proof}[证明]
我们有
\[
\widehat{\nabla}_{\hbar\partial_\hbar}(s_i)=\hbar\partial_\hbar s_i-\hbar^{-1}E*s_i+\mathrm{Gr}_E(s_i).
\]
由于 $0=\nabla_E s_i+\hbar^{-1}E*s_i$，可将其改写为
\[
\widehat{\nabla}_{\hbar\partial_\hbar}(s_i)=\hbar\partial_\hbar s_i+\nabla_E s_i+\mathrm{Gr}_E(s_i).
\]
若用 $\beta\cup\bullet$ 记映射
\[
\alpha\mapsto\beta\cup\alpha,
\]
则可写
\[
S=S'\circ(e^{-(\sum_{k=0}^p y_kT_k)/\hbar}\cup\bullet)
\]
其中 $S'$ 是适当的 $\cO_{\widetilde{\cM}}\{\hbar,\hbar^{-1}\}$-线性映射，这里用到命题 2.15 的公式。定义一个元素
\[
s\in H^*(X,\C)\otimes_\C\cO_{\widetilde{\cM}}\{\hbar,\hbar^{-1}\}
\]
为 $d$ 次齐次的，如果
\[
(\hbar\partial_\hbar+\nabla_E+\mathrm{Gr}_E)(s)=ds.
\]
于是注意，形如
\[
\hbar^{-(n+1)}T^j\kappa_1^{\int_\beta T_1}\cdots\kappa_p^{\int_\beta T_p}y_{p+1}^{n_{p+1}}\cdots y_m^{n_m}
\]
的表达式是齐次的，其次数为
\[
-(n+1)+((2\dim_\C X-\deg T_j)/2-1)+\int_\beta c_1(\cT_X)+\sum_{k=p+1}^m n_k(1-(\deg T_k)/2).
\]
另一方面，注意
\[
\langle\psi^nT_i,T_j,T_{p+1}^{n_{p+1}},\dots,T_m^{n_m}\rangle_{0,\beta}
\]
仅当（由式 (2.2)）
\[
2+\sum_{k=p+1}^m n_k+\dim_\C X-3+\int_\beta c_1(\cT_X)=n+\frac{\deg T_i+\deg T_j}{2}+\sum_{k=p+1}^m\frac{n_k\deg T_k}{2}
\]
时才非零。这说明事实上 $S'(T_i)$ 是 $(\deg T_i)/2-1$ 次齐次的，因此 $S'$ 不改变次数。于是
\[
(\hbar\partial_\hbar+\nabla_E+\mathrm{Gr}_E)\circ S'=S'\circ(\hbar\partial_\hbar+\nabla_E+\mathrm{Gr}_E),
\]
从而
\[
(\hbar\partial_\hbar+\nabla_E+\mathrm{Gr}_E)s_i=S'\big((\hbar\partial_\hbar+\nabla_E+\mathrm{Gr}_E)(e^{-(\sum_{k=0}^p y_kT_k)/\hbar}\cup T_i)\big).
\]

% [原文 p.53]
另一方面，
\begin{align*}
&(\hbar\partial_\hbar+E+\mathrm{Gr}_E)(e^{-y_0/\hbar}(e^{-(\sum_{k=1}^p y_kT_k)/\hbar}\cup T_i))\\
&\quad=(\hbar\partial_\hbar(e^{-y_0/\hbar})+E(e^{-y_0/\hbar}))(e^{-(\sum_{k=1}^p y_kT_k)/\hbar}\cup T_i)\\
&\qquad+e^{-y_0/\hbar}(\hbar\partial_\hbar+E+\mathrm{Gr}_E)(e^{-(\sum_{k=1}^p y_kT_k)/\hbar}\cup T_i)\\
&\quad=\Big(\frac{y_0}{\hbar}e^{-y_0/\hbar}-\frac{y_0}{\hbar}e^{-y_0/\hbar}\Big)(e^{-(\sum_{k=1}^p y_kT_k)/\hbar}\cup T_i)\\
&\qquad+e^{-y_0/\hbar}\Big(\hbar\partial_\hbar(e^{-(\sum_{k=1}^p y_kT_k)/\hbar}\cup T_i)\\
&\qquad\qquad+\sum_{j=1}^p-c_jT_j\hbar^{-1}\cup e^{-(\sum_{k=1}^p y_kT_k)/\hbar}\cup T_i\\
&\qquad\qquad+\sum_{n_1,\dots,n_p\ge0}\Big(\frac{y_1^{n_1}\cdots y_p^{n_p}}{n_1!\cdots n_p!}(-\hbar)^{-(n_1+\cdots+n_p)}T_1^{n_1}\cup\cdots\cup T_p^{n_p}\cup T_i\Big)\cdot\\
&\qquad\qquad\qquad\qquad\cdot\Big((n_1+\cdots+n_p)+\frac{\deg T_i}{2}-1\Big)\Big)\\
&\quad=e^{-y_0/\hbar}\Big(\hbar\partial_\hbar(e^{-(\sum_{k=1}^p y_kT_k)/\hbar}\cup T_i)\\
&\qquad\qquad-\hbar^{-1}c_1(\cT_X)\cup e^{-(\sum_{k=1}^p y_kT_k)/\hbar}\cup T_i\\
&\qquad\qquad+\Big(\frac{\deg T_i}{2}-1\Big)e^{-(\sum_{k=1}^p y_kT_k)/\hbar}\cup T_i\\
&\qquad\qquad-\hbar\partial_\hbar(e^{-(\sum_{k=1}^p y_kT_k)/\hbar}\cup T_i)\Big)\\
&\quad=e^{-(\sum_{k=0}^p y_kT_k)/\hbar}\cup\Big(\Big(\frac{\deg T_i}{2}-1\Big)-\hbar^{-1}c_1(\cT_X)\cup\Big)T_i.
\end{align*}
于是
\begin{equation}
(\hbar\partial_\hbar+\nabla_E+\mathrm{Gr}_E)\circ S=S\circ(\hbar\partial_\hbar+\nabla_E+\mathrm{Gr}_E-\hbar^{-1}c_1(\cT_X)\cup\bullet), \tag{2.12}
\end{equation}
这便是所需结果。
\end{proof}

% [Tutor Note: 式 (2.12) 与上式末行中的 $\cup$ 是算子作用记号（把 $T_i$ 之类送进去），不是普通杯积的右因子；“$-\hbar^{-1}c_1(\cT_X)\cup\bullet$”表示复合一个算子。作者特别提醒 $\mathrm{Gr}_E$ 与 $c_1(\cT_X)\cup\bullet$ 不交换：前者按次数线性作用，后者与上同调类做杯积，两者次序不同会差出 $c_1$ 的项，这正是下一证明中四个符号项相消的关键。]

因此诸截面 $s_i$ 在由 $\partial_{y_i}$ 给出的方向上是平坦的，但在由 $\partial_\hbar$ 给出的方向上不是。这并不奇怪，因为微分方程
\[
\widehat{\nabla}_{\hbar\partial_\hbar}s=0
\]
事实上在 $\hbar=0$ 与 $\hbar=\infty$ 处有奇点。因此，人们会预期该方程的任何解都是多值的。确实，只要允许 $\log\hbar$ 的幂级数，我们就能写下一个多值解，如下所述。

\begin{proposition}
设 $\hbar^{-\mathrm{Gr}_E}\hbar^{c_1(\cT_X)\cup\bullet}$ 为自同态
\begin{align*}
\exp(-\mathrm{Gr}_E\log\hbar)\exp((c_1(\cT_X)\cup\bullet)\log\hbar):{}&H^*(X,\C)\otimes_\C\cO_{\widetilde{\cM}}\{\hbar,\hbar^{-1}\}[\![\log\hbar]\!]\\
&\to H^*(X,\C)\otimes_\C\cO_{\widetilde{\cM}}\{\hbar,\hbar^{-1}\}[\![\log\hbar]\!].
\end{align*}
则
\[
\widehat{\nabla}_{\partial_{y_i}}S(\hbar^{-\mathrm{Gr}_E}\hbar^{c_1(\cT_X)\cup\bullet}T_i)=0
\]
并且
\[
\widehat{\nabla}_{\hbar\partial_\hbar}S(\hbar^{-\mathrm{Gr}_E}\hbar^{c_1(\cT_X)\cup\bullet}T_i)=0.
\]
\end{proposition}

% [原文 p.54]
\begin{proof}[证明]
第一个为零是因为对任意 $j$ 有 $\widehat{\nabla}_{\partial_{y_i}}S(T_j)=0$。对第二个，由式 (2.12) 有
\begin{align*}
&\widehat{\nabla}_{\hbar\partial_\hbar}\big(S(\hbar^{-\mathrm{Gr}_E}\hbar^{c_1(\cT_X)\cup\bullet}T_i)\big)\\
&\quad=S\Big((\hbar\partial_\hbar+\nabla_E+\mathrm{Gr}_E-\hbar^{-1}c_1(\cT_X)\cup\bullet)(\hbar^{-\mathrm{Gr}_E}\hbar^{c_1(\cT_X)\cup\bullet}T_i)\Big)\\
&\quad=S\Big(-\mathrm{Gr}_E\hbar^{-\mathrm{Gr}_E}\hbar^{c_1(\cT_X)\cup\bullet}T_i+\hbar^{-\mathrm{Gr}_E}\hbar^{c_1(\cT_X)\cup\bullet}(c_1(\cT_X)\cup T_i)\\
&\qquad\qquad+\mathrm{Gr}_E\hbar^{-\mathrm{Gr}_E}\hbar^{c_1(\cT_X)\cup\bullet}T_i-\hbar^{-1}\hbar^{-\mathrm{Gr}_E+1}\hbar^{c_1(\cT_X)\cup\bullet}(c_1(\cT_X)\cup T_i)\Big)\\
&\quad=0
\end{align*}
这由直接计算可得，须记住 $\mathrm{Gr}_E$ 与 $c_1(\cT_X)\cup\bullet$ 并不交换！
\end{proof}

值得更详细地探讨这意味着什么，办法是考察常微分方程 $\widehat{\nabla}_{\hbar\partial_\hbar}s=0$ 在 $\hbar=\infty$ 处的行为。用 $q=\hbar^{-1}$ 改写该方程，得到
\[
q\partial_q s=-qE*s+\mathrm{Gr}_E(s).
\]
该方程在 $q=0$ 处有\emph{正则奇点}。一般地，以 $q$ 为坐标、以原点为中心的圆盘 $\Delta\subseteq\C$ 上的形如
\begin{equation}
q\partial_q s=A(q)\cdot s, \tag{2.13}
\end{equation}
的常微分方程，其中未知函数 $s$ 是取值于向量空间 $\C^n$ 的全纯映射，$A$ 是全纯映射 $A:\Delta\to\mathrm{End}(\C^n)$，则称其在 $q=0$ 处有\emph{正则奇点}。在我们的例子中，$A(q)$ 是自同态 $-qE*\bullet+\mathrm{Gr}_E$。注意方程 $\widehat{\nabla}_{\hbar\partial_\hbar}s=0$ 在 $\hbar=0$ 处并没有正则奇点，因为它取 $\hbar\partial_\hbar s=B(\hbar)\cdot s$ 的形式，其中 $B(\hbar)=\hbar^{-1}E*\bullet-\mathrm{Gr}$ 在 $\hbar=0$ 处有极点。

具有正则奇点的常微分方程理论当然在经典上已发展得很完善；特别地，可以在 $\Delta\setminus\{0\}$ 上找到式 (2.13) 的 $n$ 个线性无关的多值收敛解。就当前情形而言，这些解必须与命题 2.19 中描述的形式解相合。这说明这些形式解是收敛的，特别地，量子微分方程的解 $s_i$ 是收敛的。

% [Tutor Note: p.54 的论证方向值得留意：不是先证明形式解收敛，而是先用经典 ODE 理论（正则奇点处存在 $n$ 个线性无关的多值收敛解）得到 $n$ 个解，再断言它们必与命题 2.19 的形式解相合——因为后者也是同一方程的解。注意方程在 $\hbar=0$ 处并非正则奇点（$B(\hbar)$ 有极点），故该论证只适用于 $\hbar=\infty$ 一侧。]

要了解更多上述语境中关于带奇点常微分方程的理论，[103] 给出了极好的导引。

由上述讨论可知，我们可以把映射 $S$ 看得不只是纯形式的，而视为 $\cO_{\widetilde{\cM}}\{\hbar,\hbar^{-1}\}$-线性映射
\[
S:H^*(X,\C)\otimes_\C\cO_{\widetilde{\cM}}\{\hbar,\hbar^{-1}\}\to H^*(X,\C)\otimes_\C\cO_{\widetilde{\cM}}\{\hbar,\hbar^{-1}\}
\]
其中
\[
S(\alpha)=\alpha-\sum_{j=0}^m\Big\langle\Big\langle\frac{\alpha}{\hbar+\psi},T_j\Big\rangle\Big\rangle T^j.
\]

\subsubsection{半无限 Hodge 结构变分}

下面我们在这一图景中引入更多的结构。我们将定义半无限 Hodge 结构变分的概念，这一概念由 Barannikov 在 [5] 中引入。就目前而言，这些结构看起来也许只是把上面研究过的种种数据编码起来的一种不必要地复杂的方式，但正如我们将看到的，它在 B-模型一侧自然地出现。

% [原文 p.55]
\begin{definition}
由空间 $\cM$ 参数化的一个\emph{半无限 Hodge 结构变分}是指一个有限秩局部自由 $\cO_\cM\{\hbar\}$-模 $\cE$，连同平坦联络
\[
\nabla:\cE\to\Omega^1_\cM\otimes\hbar^{-1}\cE
\]
与配对
\[
(\cdot,\cdot)_\cE:\cE\times\cE\to\cO_\cM\{\hbar\},
\]
满足
\begin{enumerate}[label=(\arabic*)]
\item $(s_1,s_2)_\cE(\hbar)=(s_2,s_1)_\cE(-\hbar)$。这里我们把 $(s_1,s_2)_\cE$ 视为 $\hbar$ 的函数；在右端，我们把 $\hbar$ 换成了 $-\hbar$。
\item $(f(-\hbar)s_1,s_2)_\cE=(s_1,f(\hbar)s_2)_\cE=f(\hbar)(s_1,s_2)_\cE$，对 $f(\hbar)\in\cO_\cM\{\hbar\}$。
\item $Y(s_1,s_2)_\cE=(\nabla_Y s_1,s_2)_\cE+(s_1,\nabla_Y s_2)_\cE$。
\item 该配对是\emph{非退化}的，即诱导配对
\[
(\cE/\hbar\cE)\otimes_{\cO_\cM}(\cE/\hbar\cE)\to\cO_\cM
\]
非退化。
\end{enumerate}
此半无限 Hodge 结构变分上的一个\emph{分次}是指一个 $\C$-线性自同态 $\mathrm{Gr}:\cE\to\cE$，使得存在 $\cM$ 上的向量场 $E$ 与常数 $D\in\C$，满足
\begin{enumerate}[label=(\arabic*),start=5]
\item $\mathrm{Gr}(fs)=((\hbar\partial_\hbar+E)f)s+f\,\mathrm{Gr}(s)$，对 $f\in\cO_\cM\{\hbar\}$，$s\in\cE$。
\item $[\mathrm{Gr},\nabla_Y]=\nabla_{[E,Y]}$，对 $\cM$ 上所有向量场 $Y$。
\item $(\hbar\partial_\hbar+E)(s_1,s_2)_\cE=(\mathrm{Gr}(s_1),s_2)_\cE+(s_1,\mathrm{Gr}(s_2))_\cE+D(s_1,s_2)_\cE$。
\end{enumerate}
\end{definition}

\begin{example}
量子上同调给出半无限 Hodge 结构变分的一个例子。照例取 $T_0,\dots,T_m$ 为 $H^*(X,\C)$ 的一组基，其中 $X$ 只有偶上同调，我们取
\[
\cE=H^*(X,\C)\otimes_\C\cO_{\widetilde{\cM}}\{\hbar\}.
\]
取平坦联络为 $\widehat{\nabla}$，即
\[
\widehat{\nabla}_X(Y)=\nabla_X(Y)+\hbar^{-1}X*Y,
\]
其中 $\nabla$ 表示 $\cE$ 上的标准平坦联络。配对为
\[
(s_1,s_2)_\cE=\int_X s_1(-\hbar)\cup s_2(\hbar).
\]
条件 (1) 与 (2) 由该定义显然成立。对 (3)，注意对 $\widetilde{\cM}$ 上的向量场 $Y$，
\begin{align*}
Y(s_1,s_2)_\cE &= \int_X Y(s_1(-\hbar)\cup s_2(\hbar))\\
&= \int_X (Ys_1(-\hbar))\cup s_2(\hbar)+s_1(-\hbar)\cup(Ys_2(\hbar))
\end{align*}
而
\begin{align*}
&(\widehat{\nabla}_Y s_1,s_2)_\cE+(s_1,\widehat{\nabla}_Y s_2)_\cE\\
&\quad=\int_X (Ys_1(-\hbar)-\hbar^{-1}Y*s_1(-\hbar))\cup s_2(\hbar)+s_1(-\hbar)\cup(Ys_2(\hbar)+\hbar^{-1}Y*s_2(\hbar))\\
&\quad=\int_X Ys_1(-\hbar)\cup s_2(\hbar)+s_1(-\hbar)\cup(Ys_2(\hbar))
\end{align*}

% [原文 p.56]
这里用到
\[
\int_X(Y*Z)\cup W=\int_X Z\cup(Y*W)
\]
（由式 (2.4)）。

(4) 显然，因为 $\cE/\hbar\cE\cong H^*(X,\C)\otimes_\C\cO_{\widetilde{\cM}}$，而诱导配对就是通常的杯积。

分次算子为
\[
\mathrm{Gr}=\hbar\partial_\hbar+E+\mathrm{Gr}_E,
\]
其中 $E$ 是例 2.10 中 $\widetilde{\cM}$ 上的 Euler 向量场，且照常，$\mathrm{Gr}_E$ 是 $\cO_{\widetilde{\cM}}\{\hbar\}$-线性的，满足
\[
\mathrm{Gr}_E(T_i)=\Big(\frac{\deg T_i}{2}-1\Big)T_i.
\]
注意该算子已经出现在命题 2.18 的证明中。我们验证条件 (5)–(7)。首先，(5) 由 $\hbar\partial_\hbar$ 与 $\nabla_E$ 的 Leibniz 律显然成立。对 (6)，首先观察到：若 (6) 对向量场 $Y$ 成立，则对 $fY$（$f\in\cO_{\widetilde{\cM}}$）也成立。又若左、右两端作用在 $\cE$ 的截面 $s$ 上相等，则作用在 $\cE$ 的截面 $fs$ 上也相等，其中 $f\in\cO_{\widetilde{\cM}}\{\hbar\}$。因此只要取 $Y=\partial_{y_i}$、作用在截面 $T_j$ 上来验证 (6) 即可。左端为
\begin{align*}
[\mathrm{Gr},\widehat{\nabla}_{\partial_{y_i}}](T_j)
&= \mathrm{Gr}(\widehat{\nabla}_{\partial_{y_i}}T_j)-\widehat{\nabla}_{\partial_{y_i}}(\mathrm{Gr}(T_j))\\
&= \mathrm{Gr}(\hbar^{-1}T_i*T_j)-\Big(\frac{\deg T_j}{2}-1\Big)\hbar^{-1}T_i*T_j\\
&= \hbar^{-1}\Big(-T_i*T_j+\mathrm{Gr}(T_i*T_j)-\Big(\frac{\deg T_j}{2}-1\Big)T_i*T_j\Big)\\
&= \hbar^{-1}\Big(\mathrm{Gr}(T_i*T_j)-\frac{\deg T_j}{2}T_i*T_j\Big).
\end{align*}
现在，若把 $T_i*T_j$ 视为 $\widetilde{\cM}$ 上的向量场——这是在把 $\widetilde{\cM}\times H^*(X,\C)$ 与 $\widetilde{\cM}$ 的切丛等同之后——那么我们知道 $T_i*T_j$ 关于 Euler 向量场 $E$ 是 $(\deg T_i+\deg T_j)/2-1$ 次齐次的，这一点由例 2.10 的讨论可得。于是
\[
\mathrm{Gr}(T_i*T_j)=\Big(\frac{\deg T_i+\deg T_j}{2}-1\Big)T_i*T_j.
\]
因此 (6) 的左端变为
\[
\hbar^{-1}\Big(\frac{\deg T_i}{2}-1\Big)(T_i*T_j).
\]
(6) 的右端为
\[
\widehat{\nabla}_{[E,\partial_{y_i}]}T_j=\hbar^{-1}\Big(\frac{\deg T_i}{2}-1\Big)T_i*T_j,
\]
这正是所需的。

% [原文 p.57]
最后，对 (7)，
\begin{align*}
&(\hbar\partial_\hbar+E)(s_1,s_2)_\cE\\
&\quad=\int_X (\hbar\partial_\hbar+E)(s_1(-\hbar))\cup s_2(\hbar)+s_1(-\hbar)\cup(\hbar\partial_\hbar+E)(s_2(\hbar))\\
&\quad=(\mathrm{Gr}(s_1),s_2)_\cE+(s_1,\mathrm{Gr}(s_2))_\cE-\int_X \mathrm{Gr}_E(s_1(-\hbar))\cup s_2(\hbar)+s_1(-\hbar)\cup\mathrm{Gr}_E(s_2(\hbar))\\
&\quad=(\mathrm{Gr}(s_1),s_2)_\cE+(s_1,\mathrm{Gr}(s_2))_\cE+(2-\dim_\C X)(s_1,s_2)_\cE.\qedhere
\end{align*}
\end{example}

% [Tutor Note: 例 2.21 在原文中跨越 pp.55–57，结尾带证毕符号 □（见 p.57 图），故本译文的 example 环境也延续到该证明结束，并在末行用 \qedhere 复现方框。另外 (4) 的“显然”用到 $\cE/\hbar\cE\cong H^*(X,\C)\otimes_\C\cO_{\widetilde{\cM}}$ 这一同构（$T_0,\dots,T_m$ 为 $H^*(X,\C)$ 的基）。]

\subsubsection{半无限 Hodge 结构变分的移动子空间实现}

我们可以用单一空间内的一个移动子空间来重新表述这一概念；这更让人联想到 Hodge 结构变分的通常概念。假设给定复流形 $\cM$ 上的一个半无限 Hodge 结构变分 $(\cE,\nabla)$，并假设 $\cM$ 单连通。考虑空间
\begin{equation}
\cH=\{s\in\Gamma(\cM,\cE\otimes_{\cO_\cM\{\hbar\}}\cO_\cM\{\hbar,\hbar^{-1}\})\mid\nabla s=0\}. \tag{2.14}
\end{equation}
由于 $\nabla$ 平坦，事实上可以看出 $\cH$ 是秩与 $\cE$ 相同的自由 $\C\{\hbar,\hbar^{-1}\}$-模。

$\cE$ 上的配对 $(\cdot,\cdot)_\cE$ 也在 $\cH$ 上定义一个辛形式，由
\begin{equation}
\overline{\Omega}(s_1,s_2):=\mathrm{Res}_{\hbar=0}(s_1,s_2)_\cE\,\dd\hbar, \tag{2.15}
\end{equation}
给出，即 $\overline{\Omega}(s_1,s_2)$ 是 $(s_1,s_2)_\cE$ 中 $\hbar^{-1}$ 的系数。注意这是 $\C$ 中的常数，因为对 $\cM$ 上任一向量场 $Y$，
\[
Y(s_1,s_2)_\cE=(\nabla_Y s_1,s_2)_\cE+(s_1,\nabla_Y s_2)_\cE=0,
\]
这里用到定义 2.20 的 (3) 以及 $s_1,s_2$ 的平坦性。

现在，对每一点 $x\in\cM$，我们得到 $\cE_x$ 到 $\cH$ 的一个嵌入：把 $s\in\cE_x$ 送到满足 $s'(x)=s$ 的截面 $s'\in\cH$；换言之，我们把 $s$ 延拓为 $\cE\otimes_{\cO_\cM\{\hbar\}}\cO_\cM\{\hbar,\hbar^{-1}\}$ 的平坦截面。于是我们得到 $\cH$ 的一族由 $\cM$ 参数化的变动子空间 $\cE_x$。

我们可以借助 $\cH$ 的半无限子空间的\emph{Pressley–Segal Grassmann 流形}来理解这些变动子空间。假设在上述半无限 Hodge 结构变分中，$\cE$ 是秩为 $N$ 的自由 $\cO_\cM\{\hbar\}$-模。则 $\cH$ 的 Pressley–Segal Grassmann 流形写作
\[
L\,\mathrm{GL}_N(\C)/L^+\mathrm{GL}_N(\C).
\]
这里 $L\,\mathrm{GL}_N(\C)$ 是 $\mathrm{GL}_N(\C)$ 的环路群，由光滑映射 $\gamma:S^1\to\mathrm{GL}_N(\C)$ 组成，其中我们把 $S^1\subseteq\C$ 视为以原点为中心、半径为某个 $\epsilon>0$ 的圆。于是我们有子群 $L^+\mathrm{GL}_N(\C)$，由那些作为全纯映射
\[
\{\hbar\mid|\hbar|<\epsilon\}\to\mathrm{GL}_N(\C)
\]
的边界值的映射组成。

这一构造在我们情形中的意义如下。假设我们有 $\C^N\otimes_\C\C\{\hbar,\hbar^{-1}\}$ 的一个秩为 $N$ 的 $\C\{\hbar\}$-子模。那么我们可以为该子模选取一组 $\C\{\hbar\}$-基 $s_1(\hbar),\dots,s_N(\hbar)$；等价地，选取一个以 $\C\{\hbar,\hbar^{-1}\}$ 的元素为元素的 $N\times N$ 可逆矩阵，使其各列张成该 $\C\{\hbar\}$-子模。现在，选取 $\epsilon>0$ 使该矩阵的所有元素在 $\{\hbar\mid0<|\hbar|<2\epsilon\}$ 上收敛，我们便可把每个元素限制到半径为
% [原文 p.58]
$\epsilon$ 的圆上，得到 $L\,\mathrm{GL}_N(\C)$ 的一个元素。$L^+\mathrm{GL}_N(\C)$ 的元素是以 $\C\{\hbar\}$ 中元素为元素的矩阵，因而其作用只是改变该 $\C\{\hbar\}$-子模的基而不改变子模。所以特别地，$\cH$ 的每个秩为 $N$ 的 $\C\{\hbar\}$-子模都由 Pressley–Segal Grassmann 流形中的一个点代表。于是半无限 Hodge 结构变分 $\cE$ 诱导出映射
\[
\cM\to L\,\mathrm{GL}_N(\C)/L^+\mathrm{GL}_N(\C).
\]
该映射可更具体地描述如下。假设 $e_1,\dots,e_N$ 是 $\cE$ 的截面，使得对每一点 $x\in\cM$，$e_1(x),\dots,e_N(x)$ 下降到 $\cE_x/\hbar\cE_x$ 的一组 $\C$-基。另一方面，设 $s_1,\dots,s_N$ 是 $\cH$ 的一组基。则存在矩阵 $M=(M_{ij})_{1\le i,j\le N}$ 使得
\[
s_i=\sum_{j=1}^N M_{ij}e_j,
\]
其中 $M_{ij}\in\cO_\cM\{\hbar,\hbar^{-1}\}$。在点 $x_0$ 处，自然有
\[
s_i(x_0)=\sum_{j=1}^N M_{ij}(x_0)e_j(x_0),
\]
而若 $M^{-1}=(M^{ij})_{1\le i,j\le N}$，则
\[
e_i(x_0)=\sum_{j=1}^N M^{ij}(x_0)s_j(x_0).
\]
于是 $\cE_{x_0}$ 通过映射
\[
e_i(x_0)\mapsto\sum_{j=1}^N M^{ij}(x_0)s_j
\]
嵌入到 $\cH$ 中。因此，在 $\cH$ 的基 $s_1,\dots,s_N$ 下，$M^{-1}$ 的各列生成 $\cE_{x_0}$ 的像。于是映射 $\cM\to L\,\mathrm{GL}_N(\C)/L^+\mathrm{GL}_N(\C)$ 可显式地描述为 $x\mapsto M^{-1}(x)$。

用 $M$ 来描述的好处在于，它在由量子上同调产生的情形中完全说得通：那里我们有空间 $\widetilde{\cM}$，由于变量 $y_0,y_{p+1},\dots,y_m$ 是形式的，它并不完全是复流形。此时，$\cH$ 中的变动子空间族 $\cE_x$ 可看作由 $M^{-1}$ 给出，即一个以 $\cO_{\widetilde{\cM}}\{\hbar,\hbar^{-1}\}$ 的元素为元素的 $N\times N$ 矩阵。

% [Tutor Note: p.57–58 的对应关系是：秩 $N$ 的 $\C\{\hbar\}$-子模 $\leftrightarrow$ 一个 $N\times N$ 可逆矩阵的列空间 $\leftrightarrow$ Pressley–Segal Grassmann 流形 $L\mathrm{GL}_N(\C)/L^+\mathrm{GL}_N(\C)$ 中的一点。技术要点在“先取 $\epsilon>0$ 使元素在 $0<|\hbar|<2\epsilon$ 收敛，再限制到 $|\hbar|=\epsilon$ 的圆”——晚一步取圆环才会有 $L^+\mathrm{GL}_N(\C)$ 的等价作用。]

\begin{example}
回到由量子上同调给出的例 2.21，我们取基 $e_1,\dots,e_N$ 为 $T_0,\dots,T_m$。我们还有 §2.1.4 的平坦截面 $s_0,\dots,s_m$。矩阵 $M$ 恰好由函数 $S$ 给出，因此相应半无限 Hodge 结构变分的移动子空间实现由 $S^{-1}$ 定义。我们定义
\[
\mathbb{J}=S^{-1}:H^*(X,\C)\otimes_\C\cO_{\widetilde{\cM}}\{\hbar,\hbar^{-1}\}\to H^*(X,\C)\otimes_\C\cO_{\widetilde{\cM}}\{\hbar,\hbar^{-1}\}.
\]
由于用 $H^*(X,\C)$ 的基 $\{T_i\}$ 定义的 $S$ 的矩阵元素属于 $\cO(\PP^1\setminus\{0\})$，并且在 $\PP^1\setminus\{0\}$ 的每一点可逆，$\mathbb{J}$ 的矩阵元素也属于 $\cO(\PP^1\setminus\{0\})$。
\end{example}

% [原文 p.59]
\begin{proposition}
对 $\alpha\in H^*(X,\C)$，
\begin{align*}
\mathbb{J}(\alpha) &= e^{(\sum_{k=0}^p y_kT_k)/\hbar}\cup\Big(\alpha+\sum_{\beta,n_k}\sum_{i=0}^m\Big\langle\alpha,T_{p+1}^{n_{p+1}},\dots,T_m^{n_m},\frac{T_i}{\hbar-\psi}\Big\rangle_{0,\beta}T^i\\
&\qquad\qquad\qquad\qquad\cdot\Big(\prod_{k=1}^p e^{y_k\int_\beta T_k}\Big)\frac{y_{p+1}^{n_{p+1}}\cdots y_m^{n_m}}{n_{p+1}!\cdots n_m!}\Big).
\end{align*}
\end{proposition}

\begin{proof}[证明]
我们首先证明：取 $\cE=H^*(X,\C)\otimes_\C\cO_{\widetilde{\cM}}\{\hbar\}$ 如例 2.21，对 $\alpha,\beta\in H^*(X,\C)$ 有
\[
(S(\alpha),S(\beta))_\cE=(\alpha,\beta)_\cE.
\]
确实，应用定义 2.20 的性质 (3)，可见
\begin{align*}
\partial_{y_i}(S(\alpha),S(\beta))_\cE &= (\widehat{\nabla}_{\partial_{y_i}}S(\alpha),S(\beta))_\cE+(S(\alpha),\widehat{\nabla}_{\partial_{y_i}}S(\beta))_\cE\\
&= 0,
\end{align*}
因为 $S(\alpha),S(\beta)$ 满足量子微分方程。于是 $(S(\alpha),S(\beta))_\cE$ 是 $\widetilde{\cM}$ 上的常函数。因此我们可以试着通过令 $y_0=y_{p+1}=\cdots=y_m=0$ 并让 $y_i\to-\infty$（即对 $1\le i\le p$ 令 $\kappa_i\to0$）来计算它。取该极限，再利用命题 2.15 的公式，我们得到
\begin{align*}
(S(\alpha),S(\beta))_\cE &= \int_X \big(e^{(\sum_{k=0}^p y_kT_k)/\hbar}\cup\alpha\big)\cup\big(e^{-(\sum_{k=0}^p y_kT_k)/\hbar}\cup\beta\big)\\
&= \int_X \alpha\cup\beta.
\end{align*}
这表明 $S^{-1}$ 是 $S$ 关于内积 $(\cdot,\cdot)_\cE$ 的伴随。于是我们可以用
\begin{align*}
(T_i,e^{-(\sum_{k=0}^p y_kT_k)/\hbar}\cup\mathbb{J}(T_j))_\cE &= (e^{(\sum_{k=0}^p y_kT_k)/\hbar}\cup T_i,\mathbb{J}(T_j))_\cE\\
&= (S(e^{(\sum_{k=0}^p y_kT_k)/\hbar}\cup T_i),T_j)_\cE
\end{align*}
来计算 $\mathbb{J}$，于是把 $e^{-(\sum_{k=0}^p y_kT_k)/\hbar}\cup\mathbb{J}(T_j)$ 中 $T^i$ 的系数读作 $S(e^{(\sum_{k=0}^p y_kT_k)/\hbar}\cup T_i)(-\hbar)$ 中 $T^j$ 的系数，这里用到命题 2.15 的公式。
\end{proof}

一种看待 $\mathbb{J}$ 的方式是：它给出丛的嵌入
\[
\mathbb{J}:\cE=H^*(X,\C)\otimes_\C\cO_{\widetilde{\cM}}\{\hbar\}\to\cH_{\widetilde{\cM}}=H^*(X,\C)\otimes_\C\cO_{\widetilde{\cM}}\{\hbar,\hbar^{-1}\},
\]
把 $\cE$ 的平坦截面送到 $\cH_{\widetilde{\cM}}$ 的常截面。

\subsubsection{从半无限 Hodge 结构变分到 Frobenius 流形}

下面我们演示 Barannikov 的、从半无限 Hodge 结构变分得到 Frobenius 流形的技巧。设 $(\cM,\cE,\nabla)$ 是一个半无限 Hodge 结构变分，其中 $\cM$ 是单连通复流形，或是复流形沿单连通子流形的完备化；固定基点 $0\in\cM$。这确定了如式 (2.14) 中的空间 $\cH$，以及一个把 $\cE_0$ 作为 $\C\{\hbar\}$-子模包含进 $\cH$ 的包含 $\cE_0\hookrightarrow\cH$。

一个\emph{对置子空间}是 $\cH$ 的一个 $\cO(\PP^1\setminus\{0\})$-子模 $\cH_-$，使得自然映射
\[
\cH_-\oplus\cE_0\to\cH
\]
是同构。

% [原文 p.60]
首先注意，投影 $\cE_0\cap\hbar\cH_-\to\cE_0/\hbar\cE_0$ 是同构。确实，该映射是单射，因为其核为 $\hbar\cE_0\cap\hbar\cH_-=\hbar(\cE_0\cap\cH_-)=0$。为证满射，若 $s\in\cE_0$，则 $\hbar^{-1}s=s'+s''$，其中 $s'\in\cH_-$，$s''\in\cE_0$。于是 $\hbar s'\in\cE_0\cap\hbar\cH_-$ 且 $\hbar s'\equiv s\bmod\hbar\cE_0$。类似地，$\cE_0\cap\hbar\cH_-\to\hbar\cH_-/\cH_-$ 也是同构。

这些同构于是给出同构
\[
\cE_0\cong(\cE_0\cap\hbar\cH_-)\otimes_\C\C\{\hbar\}\cong\Big(\frac{\cE_0}{\hbar\cE_0}\Big)\otimes_\C\C\{\hbar\}\cong\Big(\frac{\hbar\cH_-}{\cH_-}\Big)\otimes_\C\C\{\hbar\}.
\]
确实，映射 $(\cE_0\cap\hbar\cH)\otimes_\C\C\{\hbar\}\to\cE_0$ 是显然的那个，$s\otimes f\mapsto s\cdot f$。

% [Tutor Note（原文笔误，照译未改）: p.60 上式中的 $(\cE_0\cap\hbar\cH)\otimes_\C\C\{\hbar\}\to\cE_0$ 应为 $(\cE_0\cap\hbar\cH_-)\otimes_\C\C\{\hbar\}\to\cE_0$。依据：同页上一行的同构链写的是 $\cE_0\cong(\cE_0\cap\hbar\cH_-)\otimes_\C\C\{\hbar\}$，且本段开头证明的就是 $\cE_0\cap\hbar\cH_-\to\cE_0/\hbar\cE_0$ 为同构。]

由 §2.1.6，我们可以把 $\cE$ 看成 $\cM$ 上平凡丛 $\cH\times\cM$ 的子丛。一般地，当我们有一个向量空间 $V$ 时，我们用 $V_\cM$ 记以 $V$ 为纤维的平凡向量丛，而元素 $s\in V$ 给出 $V_\cM$ 的一个常截面，我们也用 $s$ 记之。于是我们有向量丛 $\cH_\cM$，它包含 $\cE$ 作为子丛，也包含平凡子丛 $\cH_{-,\cM}$。投影 $\cE\cap\hbar\cH_{-,\cM}\to(\hbar\cH_-/\cH_-)_\cM$ 在 $0$ 处是同构，因而在 $0$ 的某个开邻域内是同构。把 $\cM$ 换为该开邻域，我们得到平凡化
\begin{equation}
\tau:\Big(\frac{\hbar\cH_-}{\cH_-}\Big)\otimes_\C\cO_\cM\{\hbar\}\to\cE \tag{2.16}
\end{equation}
它由下式给出：对 $s\in\hbar\cH_-$，写 $s=s'+s''$，其中 $s'$ 是 $\cE$ 的截面，$s''$ 是 $\cH_{-,\cM}$ 的截面，则
\[
\tau((s\bmod\cH_-)\otimes f)=s'\cdot f.
\]

\begin{proposition}
给定一个对置子空间 $\cH_-$ 的选取，我们有：
\begin{enumerate}[label=(\arabic*)]
\item 经由平凡化 (2.16)，$\cE$ 上的联络 $\nabla$ 给出 $(\hbar\cH_-/\cH_-)\otimes_\C\cO_\cM\{\hbar\}$ 上的一个联络 $\nabla$。我们可以把该联络写成
\[
\nabla=\dd+\hbar^{-1}A,
\]
其中 $A$ 是 $\cM$ 上取值为 $\mathrm{End}(\hbar\cH_-/\cH_-)$ 的 1-形式。对 $\cM$ 上的向量场 $X$，我们用 $A_X$ 记平凡丛 $\mathrm{End}(\hbar\cH_-/\cH_-)_\cM$ 的相应截面。
\item 若 $\cH_-$ 关于式 (2.15) 中定义的 $\overline{\Omega}$ 是各向同性的，则我们得到向量空间 $\hbar\cH_-/\cH_-$ 上的对称配对，由
\[
(s_1,s_2)_{\hbar\cH_-/\cH_-}=(s_1,s_2)_\cE|_{\hbar=\infty}
\]
给出。（这里右端我们把 $s_1,s_2\in\hbar\cH_-$ 视为 $\cE\otimes_{\cO_\cM\{\hbar\}}\cO_\cM\{\hbar,\hbar^{-1}\}$ 的平坦截面，$\cE$ 上的配对在其上线性延拓。）此外，该配对非退化，并且
\begin{equation}
(A_Xs_1,s_2)_{\hbar\cH_-/\cH_-}=(s_1,A_Xs_2)_{\hbar\cH_-/\cH_-}. \tag{2.17}
\end{equation}
\item 若该半无限 Hodge 结构变分有一个保持 $\cH_-$ 的分次算子 $\mathrm{Gr}$，则关于平凡化 (2.16)，可写
\[
\mathrm{Gr}=\hbar\partial_\hbar+E+\mathrm{Gr}_0,
\]
其中 $\mathrm{Gr}_0\in\mathrm{End}(\hbar\cH_-/\cH_-)$ 由下式定义：对 $s\in\hbar\cH_-$，
\[
\mathrm{Gr}_0(s\bmod\cH_-)=\mathrm{Gr}(s)\bmod\cH_-.
\]
此外，对 $s_1,s_2\in\hbar\cH_-/\cH_-$，
\[
-D(s_1,s_2)_{\hbar\cH_-/\cH_-}=(\mathrm{Gr}_0(s_1),s_2)_{\hbar\cH_-/\cH_-}+(s_1,\mathrm{Gr}_0(s_2))_{\hbar\cH_-/\cH_-}.
\]
\end{enumerate}
\end{proposition}

% [原文 p.61]
\begin{proof}[证明]
(1) 我们来计算 $\nabla_X(\tau(s\otimes1))$，其中 $s\in\hbar\cH_-$ 代表 $(\hbar\cH_-/\cH_-)_\cM$ 的一个常截面。作为截面，可写 $s=s'+s''$，其中 $s'\in\cE$，$s''\in\cH_{-,\cM}$，从而 $\tau(s\otimes1)=s'=s-s''$。注意在嵌入 $\cE\subseteq\cH_\cM$ 之下，$\cH_\cM$ 上的平凡联络诱导出定义该半无限 Hodge 结构变分的 $\cE$ 上的联络 $\nabla$，因为 $\cH_\cM$ 的常截面正是 $\cE\otimes_{\cO_\cM\{\hbar\}}\cO_\cM\{\hbar,\hbar^{-1}\}$ 的平坦截面。于是
\[
\hbar\nabla_X(\tau(s\otimes1))=\hbar\nabla_X(s-s'')=-\hbar X(s''),
\]
而由于 $X(s'')\in\cH_{-,\cM}$，我们有 $\hbar X(s'')\in\hbar\cH_{-,\cM}$。另一方面，
\[
\hbar\nabla_X(\tau(s\otimes1))\in\cE,
\]
所以
\[
\hbar X(s'')\in\cE\cap\hbar\cH_{-,\cM}\cong(\hbar\cH_-/\cH_-)_\cM.
\]
因此 $\hbar\nabla_X$ 定义了平凡丛 $\mathrm{End}(\hbar\cH_-/\cH_-)_\cM$ 的一个截面 $s\mapsto-\hbar X(s'')$，而 $A$ 是取值于该丛的 1-形式，由 $A_X(s)=-\hbar X(s'')$ 给出。

(2) 首先注意，$\hbar\cH_-/\cH_-$ 上给出的配对是有意义的。确实，给定 $s_1,s_2\in\hbar\cH_-/\cH_-$，令 $s_i'=\tau(s_i\otimes1)$ 为 $\cE$ 的相应截面，并照常写 $s_i=s_i'+s_i''$。由于 $s_i'$ 是 $\cE$ 的截面，$(s_1',s_2')_\cE\in\cO_\cM\{\hbar\}$。另一方面，若 $t_1,t_2$ 是 $\hbar\cH_{-,\cM}$ 的截面，则对任意 $a_1,a_2>0$，可写 $t_i=\hbar^{a_i}t_i'$，其中 $t_i'$ 是 $\cH_-$ 的截面；于是利用 $\cH_-$ 关于 $\overline{\Omega}$ 的各向同性，我们有
\[
0=\overline{\Omega}(t_1',t_2')=\pm\mathrm{Res}_{\hbar=0}\hbar^{-a_1-a_2}(t_1,t_2)_\cE\,\dd\hbar,
\]
由此可知 $(t_1,t_2)_\cE$ 只含 $\hbar$ 的 $\le0$ 阶项。类似地，若 $t_1\in\hbar\cH_{-,\cM}$，$t_2\in\cH_{-,\cM}$，则 $(t_1,t_2)_\cE$ 只含 $\hbar$ 的 $\le-1$ 阶项。于是 $(s_1',s_2')_\cE=(s_1-s_1'',s_2-s_2'')_\cE$ 事实上是 $\cO_\cM$ 的截面，而 $(s_1',s_2')_\cE=(s_1',s_2')_\cE|_{\hbar=\infty}$ 有意义，且与 $s_1,s_2$ 在 $\hbar\cH_-$ 中代表元的选取无关。此外，这与 $(s_1,s_2)_\cE|_{\hbar=\infty}$ 相同，因为后一表达式中 $\hbar$ 的非负幂次的项与 $(s_1',s_2')_\cE$ 中对应的项一致。最后注意，由于 $s_1,s_2$ 是平坦截面，由定义 2.20 的 (3) 可知 $(s_1,s_2)_\cE$ 事实上是 $\cM$ 上的常数，因此 $(s_1,s_2)_\cE|_{\hbar=\infty}$ 给出 $\C$ 的一个确定元素。

特别地，再次由定义 2.20 的 (3)，
\begin{align*}
0 &= \hbar Y(s_1',s_2')_\cE=-(\hbar\nabla_Y s_1',s_2')_\cE+(s_1',\hbar\nabla_Y s_2')_\cE\\
&= -(A_Ys_1,s_2)_{\hbar\cH_-/\cH_-}+(s_1,A_Ys_2)_{\hbar\cH_-/\cH_-},
\end{align*}
由 (1) 即得所需公式。

$\hbar\cH_-/\cH_-$ 上配对的非退化性由定义 2.20 的 (4) 的非退化性推出。

对 (3)，首先注意由定义 2.20 的 (6)，$\mathrm{Gr}$ 把 $\cE$ 的平坦截面变为 $\cE$ 的平坦截面。因此 $\mathrm{Gr}$ 自然地作用在 $\cH$ 上。此外，由定义 2.20 的 (5)，若 $\mathrm{Gr}$ 保持 $\cH_-$，则它也保持 $\hbar\cH_-$。于是 $\mathrm{Gr}$ 定义了 $\hbar\cH_-/\cH_-$ 的一个确定的自同态 $\mathrm{Gr}_0$。$\mathrm{Gr}$ 用 $\mathrm{Gr}_0$ 的描述随即由定义 2.20 的 (5) 得到：
\begin{align*}
\mathrm{Gr}(\tau(s\otimes f)) &= \mathrm{Gr}(f\cdot s')=((\hbar\partial_\hbar+E)f)s'+f\,\mathrm{Gr}(s')\\
&= \tau(s\otimes(\hbar\partial_\hbar+E)f+\mathrm{Gr}_0(s)\otimes f).
\end{align*}
最后一个断言直接由定义 2.20 的 (7) 得到。
\end{proof}

% [Tutor Note（原文笔误，照译未改）: p.61 第 2 段“it follow from Definition 2.20, (3)”应为“it follows from ...”。另注意同段要求 $a_1,a_2>0$ 任意，故“$(t_1,t_2)_\cE$ 只含 $\le0$ 阶项”这一估计对一切足够大的 $a_1+a_2$ 成立，这正是把 $(s_1',s_2')_\cE$ 判为 $\cO_\cM$ 截面的依据。]
